\section{Conclusion}
\label{sec:conclusion}

We asked whether reduced N3 sleep shows a dose-dependent association with amyloid-beta burden in
cognitively normal adults, and whether that association is stronger than the associations with
total sleep time or sleep efficiency. The answer to both is no. Pooling 40 amyloid-PET effects
from 17 studies across 12 cohort families, N3 stage amount gives $r = +0.026$ (95\% CI $-0.153$
to $+0.203$), equivalence-tested tightly enough to exclude $r \geq 0.20$; its dose statement
converts to $+0.005$ SUVR per 10 percentage points, 1.8\% of the amyloid-negative/positive
separation; and it is weaker rather than stronger than both comparators, a gap that differential
measurement reliability makes harder to explain away rather than easier, because a single night
measures N3 more precisely than it measures either comparator.

The key takeaway is that ``slow-wave sleep'' names two different measurements with two different
evidentiary situations. The \emph{stage} metric is a tight, homogeneous null across four
independent cohort families. The \emph{spectral} metric is not refuted, but its pooled estimate
is carried entirely by one cohort family --- $r = +0.112$ overall, $-0.025$ with that family
removed --- and our spectral stratum is too imprecise ($\mathrm{MDE}\ r = 0.417$) to support any
conclusion in either direction. The confident causal framing now common in the secondary
literature rests on the second construct while the field measures and intervenes on the first.

\para{Future work.} Four steps would settle what this design cannot. An individual-participant
analysis in A4, ADNI or OASIS-3 is the only way to estimate a genuine dose-response curve, test
for thresholds, and check whether a null is masked by effect modification through APOE4, age or
baseline amyloid status. Multi-night polysomnography or home EEG in an amyloid-imaged cohort
would raise total sleep time's reliability from 0.225 toward N3's, removing the asymmetry
entirely and letting the biological comparison be made cleanly. A properly powered spectral
replication --- roughly $n \approx 190$ in a single independent cohort to detect $r = 0.20$ at
80\% power --- is needed before the Berkeley effect can be treated as established. And a
region- and sub-band-resolved analysis would test whether the medial-prefrontal, 0.6--1\,Hz
effect of \citet{mander2015} is simply averaged away by global-burden pooling.
